% Generated by experiments/scripts/export_latex.py -- do not edit by hand.
% Source: experiments/data/2-processed-tables-and-evaluations/revised-usd/same-window/eval_summary.json
\begin{table}[htbp]
\centering
\caption[빠진 계약을 0점으로 매기거나 외톨이 점으로 둘 때의 같은 기간 지표 가족 평균.]{그래프에서 빠진 계약을 0점으로 매기거나 외톨이 점(화살표가 하나도 없는 점)으로 둘 때, 같은 기간의 지표 가족 평균 켄달의 $\tau$를 보여 줘요(매칭 코호트, 계약 $n=15{,}107$개). 주 분석에서는 어떤 방법의 그래프에 없는 계약을 0점으로 매겨요. 그 계약을 외톨이 점으로 두는 것은 분석 계획에 미리 정해 둔 민감도 분석이에요. 이 규칙의 영향을 받는 계약은 엔도스랭크에서 523개, AWP에서 2{,}670개예요(80개는 두 그래프 모두에 없어요). 엔도스랭크는 고르게 다시 출발하기 때문에, 외톨이 점은 들어오는 화살표가 없는 점과 같은 점수를 받아요. AWP는 송금을 보낸 주소에서만 다시 출발해요. 그래서 외톨이 계약은 점수 0을 그대로 갖고, AWP의 계수는 두 규칙에서 똑같아요. 마지막 행은 두 점수 사이의 일치예요. 신뢰구간(CI)은 계약을 다시 뽑아 구한 95\% 백분위 부트스트랩 구간이에요(짝지은 다시 뽑기 400번, 시드 42).}
\label{tab:tie-sensitivity}
\fitwidth{%
\begin{tabular}{lcccc}
\toprule
 & \multicolumn{2}{c}{빠진 계약을 0점으로} & \multicolumn{2}{c}{빠진 계약을 외톨이 점으로} \\
\cmidrule(lr){2-3} \cmidrule(lr){4-5}
지표 가족 & EndorseRank & AWP & EndorseRank & AWP \\
\midrule
송금 & 0.158 [0.147, 0.169] & 0.572 [0.564, 0.580] & 0.158 [0.146, 0.169] & 0.572 [0.564, 0.580] \\
허락 & 0.585 [0.578, 0.592] & 0.193 [0.182, 0.203] & 0.585 [0.579, 0.592] & 0.193 [0.182, 0.203] \\
시빌 안정성 & 0.007 [-0.001, 0.015] & 0.395 [0.387, 0.403] & 0.007 [-0.002, 0.015] & 0.395 [0.387, 0.403] \\
\midrule
EndorseRank 대 AWP & \multicolumn{2}{c}{0.174 [0.163, 0.185]} & \multicolumn{2}{c}{0.174 [0.163, 0.185]} \\
\bottomrule
\end{tabular}
}
\end{table}
